% bismi ALLAH ARRAHMAN ARRAHIM




















































































\documentclass[11pt,a4paper]{article}

\usepackage[utf8]{inputenc}
\usepackage[T1]{fontenc}
\usepackage[french]{babel}
\usepackage{amsmath,amssymb,amsfonts,amsthm,mathtools}
\usepackage{booktabs}
\usepackage{geometry}
\usepackage{enumitem}
\usepackage{microtype}
\usepackage{hyperref}
\usepackage{tikz}
\usetikzlibrary{arrows.meta,positioning,calc}
\hypersetup{hidelinks}
\geometry{margin=2.5cm}

% -------------------------------------------------
% Environnements
% -------------------------------------------------
\newtheorem{definition}{Définition}[section]
\newtheorem{assumption}{Hypothèse}[section]
\newtheorem{proposition}{Proposition}[section]
\newtheorem{theorem}{Théorème}[section]
\newtheorem{corollary}{Corollaire}[section]
\newtheorem{remark}{Remarque}[section]
\newtheorem{example}{Exemple}[section]

% -------------------------------------------------
% Opérateurs
% -------------------------------------------------
\DeclareMathOperator{\prox}{prox}
\DeclareMathOperator{\argmin}{argmin}
\DeclareMathOperator{\argmax}{argmax}
\DeclareMathOperator{\dist}{dist}
\DeclareMathOperator{\conv}{conv}

% -------------------------------------------------
% Raccourcis
% -------------------------------------------------
\newcommand{\R}{\mathbb R}
\newcommand{\Hh}{\mathcal H}
\newcommand{\Mcal}{\mathcal M}
\newcommand{\Ccal}{\mathcal C}
\newcommand{\Pcal}{\mathcal P}
\newcommand{\Tcal}{\mathcal T}
\newcommand{\Rcal}{\mathcal R}
\newcommand{\Acal}{\mathcal A}
\newcommand{\Bcal}{\mathcal B}

\title{Routage variationnel PCH-ESM\\
\large Sélection de contexte par fonctions valeur et interprétation variationnelle}
\author{}
\date{}

\begin{document}

\maketitle

\begin{abstract}
Ce manuscrit formalise le routage dans le cadre PCH comme une lecture variationnelle d'une famille finie de fonctions valeur locales associées aux contextes. Le critère de décision ne repose pas sur un score appris supplémentaire : il sélectionne le contexte minimisant l'enveloppe inférieure des valeurs réduites issues de problèmes locaux convexes. On obtient ainsi un problème d'optimisation sur \(\Hh\times\Ccal\), une relaxation entropique sur le simplexe des contextes, une géométrie de cellules et de frontières, ainsi qu'une instanciation PCH-ESM où la décomposition spectrale fournit une famille finie de blocs spécialisés. Les outils classiques de Moreau, Berge, Clarke et Rockafellar--Wets sont utilisés comme instruments d'analyse ; la contribution propre est la formulation variationnelle qui fait émerger un routeur à partir de fonctions valeur.
\end{abstract}

\tableofcontents

% =================================================
\section{Introduction}
% =================================================

On considère un noyau PCH déjà défini et validé. Pour chaque entrée latente \(w\in\Hh\), plusieurs contextes admissibles \(c\in\Ccal\) sont possibles, chacun associé à un ensemble convexe non vide et fermé \(\Mcal_c\subset\Hh\) et à une fonctionnelle locale propre.

Le point conceptuel central est simple : la question \emph{quel contexte purifier ?} est formulée, dans le cadre proposé, par la comparaison de valeurs locales. Le routage n'ajoute pas de score externe ; il compare les branches d'énergie que le noyau fournit déjà.

La bascule conceptuelle est la suivante : au lieu d'apprendre un routeur séparé, on lit la décision dans une famille de fonctions valeur. Autrement dit, la sélection de contexte n'est pas imposée a posteriori ; elle émerge du minimum d'une famille finie de valeurs locales.

\begin{quote}
\emph{Le routage est l'enveloppe inférieure d'une famille finie de fonctions valeur locales.}
\end{quote}

Cette phrase ne remplace pas la théorie ; elle en constitue le point de bascule. Elle conduit à une géométrie de décision, à une stabilité par marge, à une relaxation entropique et à une compatibilité naturelle avec PCH-ESM.

\begin{remark}
Le cadre ne dépend pas du mot PCH en tant que tel : toute famille finie de fonctions valeur locales induit le même mécanisme. PCH fournit ici une instanciation naturelle, structurée et lisible de cette situation.
\end{remark}

\begin{figure}[ht]
\centering
\begin{tikzpicture}[node distance=1.7cm, >=Latex, every node/.style={align=center}]
\node (w) [draw, rounded corners, minimum width=2.8cm, minimum height=0.9cm] {$w\in\Hh$};
\node (locals) [draw, rounded corners, right=2.7cm of w, minimum width=4.2cm, minimum height=1.0cm] {Problèmes locaux\\$\widetilde J_c(\cdot;w)$};
\node (values) [draw, rounded corners, right=2.7cm of locals, minimum width=3.8cm, minimum height=1.0cm] {Valeurs réduites\\$\widetilde V_c(w)$};
\node (argmin) [draw, rounded corners, below=1.25cm of values, minimum width=3.8cm, minimum height=1.0cm] {Sélection\\$\argmin_{c\in\Ccal}\widetilde V_c(w)$};
\node (chosen) [draw, rounded corners, left=2.7cm of argmin, minimum width=3.8cm, minimum height=1.0cm] {Contexte choisi\\$c^\star(w)$};

\draw[->] (w) -- (locals);
\draw[->] (locals) -- (values);
\draw[->] (values) -- (argmin);
\draw[->] (argmin) -- (chosen);
\end{tikzpicture}
\caption{Chaîne conceptuelle du routage variationnel.}
\end{figure}

\section{Contributions}

Les contributions du manuscrit sont formulées de manière courte et vérifiable.

\begin{itemize}[leftmargin=2em]
    \item Formuler le routage comme le minimum d'une famille finie de fonctions valeur locales.
    \item Montrer que cette formulation induit une géométrie de cellules, de frontières et de marges.
    \item Faire apparaître explicitement la bascule conceptuelle entre routeur appris et lecture variationnelle.
    \item Déduire un routage souple comme relaxation entropique de la règle dure.
    \item Ajouter un résultat décisionnel fort : optimalité Bayesienne sous représentation en risque conditionnel et robustesse sous perturbation.
    \item Présenter PCH-ESM comme une instanciation naturelle de la théorie générale.
\end{itemize}

\section*{Table des notations principales}
\begin{center}
\begin{tabular}{@{}ll@{}}
\toprule
Notation & Signification \\ \midrule
\(\widetilde J_c(u;w)\) & énergie locale normalisée du contexte \(c\) \\
\(\widetilde V_c(w)\) & valeur réduite associée au contexte \(c\) \\
\(u_c^\star(w)\) & minimiseur local unique du problème normalisé \\
\(F(w)\) & valeur globale de routage, \(F(w)=\min_c \widetilde V_c(w)\) \\
\(\Acal(w)\) & ensemble des contextes actifs \\
\(\Omega_c^{\mathrm{str}}\) & cellule stricte du contexte \(c\) \\
\(\Bcal\) & ensemble des points de bord \\
\(\Rcal_\beta\) & fonctionnelle de routage souple \\
\(\Delta_K\) & simplexe des contextes \\
\bottomrule
\end{tabular}
\end{center}

% =================================================
\section{Problème local, normalisation et valeur réduite}
% =================================================

\begin{assumption}[Espace latent]
On fixe un espace euclidien réel \(\Hh=\R^d\). Les états latents purifiés vivent dans \(\Hh\).
\end{assumption}

\begin{assumption}[Famille de contextes]
On considère un ensemble fini de contextes admissibles
\[
\Ccal=\{1,\dots,K\}.
\]
À chaque contexte \(c\in\Ccal\) est associé un ensemble convexe non vide et fermé \(\Mcal_c\subset\Hh\).
\end{assumption}

\begin{assumption}[Paramètres PCH]
Pour chaque contexte \(c\), on fixe des paramètres \(\lambda_{\mathrm{ctx},c}>0\) et \(\lambda_{\mathrm{ridge},c}\ge 0\).
\end{assumption}

\begin{definition}[Énergie contextuelle]
Pour un contexte \(c\in\Ccal\), on définit
\[
J_c(u;w)=\frac12\|u-w\|^2+\lambda_{\mathrm{ctx},c}\,\dist(u,\Mcal_c)^2+\lambda_{\mathrm{ridge},c}\,\|u\|^2.
\]
\end{definition}

\begin{definition}[Opérateur PCH associé]
Pour un contexte \(c\in\Ccal\), on définit
\[
\Pcal_c(w):=\argmin_{u\in\Hh}J_c(u;w).
\]
\end{definition}

\begin{proposition}[Convexité forte et unicité]
Pour tout contexte \(c\) et toute entrée \(w\), la fonctionnelle \(u\mapsto J_c(u;w)\) est \(\mu_c\)-fortement convexe, où
\[
\mu_c:=1+2\lambda_{\mathrm{ridge},c}.
\]
En particulier, \(\Pcal_c(w)\) est un singleton.
\end{proposition}

\begin{proof}
Le terme \(u\mapsto \frac12\|u-w\|^2+\lambda_{\mathrm{ridge},c}\|u\|^2\) est \((1+2\lambda_{\mathrm{ridge},c})\)-fortement convexe. La fonction \(u\mapsto \dist(u,\Mcal_c)^2\) est convexe parce que \(\Mcal_c\) est non vide, fermé et convexe. La somme est donc fortement convexe, d'où l'existence et l'unicité du minimiseur.
\end{proof}

\begin{remark}
La normalisation par \(\mu_c\) ne change pas le minimiseur ; elle aligne les branches sur une échelle commune et isole la composante proximale du problème. C'est cette échelle qui rend la comparaison des valeurs réduites canonique.
\end{remark}

\begin{definition}[Fonction de base normalisée]
On pose
\[
f_c(u):=\lambda_{\mathrm{ctx},c}\,\dist(u,\Mcal_c)^2+\lambda_{\mathrm{ridge},c}\,\|u\|^2,
\qquad
g_c(u):=\frac{1}{\mu_c}f_c(u).
\]
Alors
\[
\widetilde J_c(u;w)=\frac{1}{\mu_c}J_c(u;w)=g_c(u)+\frac{1}{2\mu_c}\|u-w\|^2.
\]
\end{definition}

\begin{definition}[Valeur réduite]
Pour tout contexte \(c\) et toute entrée \(w\), on définit la valeur réduite
\[
\widetilde V_c(w):=\inf_{u\in\Hh}\widetilde J_c(u;w).
\]
On note aussi
\[
 u_c^\star(w):=\argmin_{u\in\Hh}\widetilde J_c(u;w),
\qquad
 e_{\mu_c g_c}(w):=\inf_{u\in\Hh}\left\{g_c(u)+\frac{1}{2\mu_c}\|u-w\|^2\right\}.
\]
\end{definition}

\begin{proposition}[Enveloppe de Moreau et gradient]
Pour tout contexte \(c\), la fonction \(g_c\) est propre, convexe et semi-continue inférieurement. On a
\[
 u_c^\star(w)=\prox_{\mu_c g_c}(w),
\qquad
 \widetilde V_c(w)=e_{\mu_c g_c}(w).
\]
De plus, \(e_{\mu_c g_c}\) est de classe \(C^1\) sur \(\Hh\) et
\[
\nabla e_{\mu_c g_c}(w)=\frac{1}{\mu_c}\bigl(w-\prox_{\mu_c g_c}(w)\bigr)
=\frac{1}{\mu_c}\bigl(w-u_c^\star(w)\bigr).
\]
Par conséquent,
\[
\nabla \widetilde V_c(w)=\frac{1}{\mu_c}\bigl(w-u_c^\star(w)\bigr).
\]
\end{proposition}

\begin{proof}
La convexité et la semi-continuité inférieure de \(g_c\) proviennent de celles de \(\dist(\cdot,\Mcal_c)^2\) et de \(\|\cdot\|^2\). Le résultat classique sur les enveloppes de Moreau d'une fonction propre, convexe et s.c.i. implique la différentiabilité \(C^1\) de \(e_{\mu_c g_c}\) et la formule du gradient. Comme \(\widetilde J_c(u;w)=g_c(u)+\frac{1}{2\mu_c}\|u-w\|^2\), les minimiseurs coïncident.
\end{proof}

\begin{remark}
Ces identités sont des résultats standards de l'analyse convexe ; elles sont rappelées seulement là où elles servent le raisonnement \cite{Moreau1965,Rockafellar1970,BauschkeCombettes2017,RockafellarWets1998}.
\end{remark}

\begin{remark}
Dans ce cadre, PCH génère naturellement une famille de valeurs locales comparables. La règle de routage consiste ensuite à prendre leur enveloppe inférieure canonique ; le mécanisme est général, et PCH en fournit une instanciation naturelle.
\end{remark}

% =================================================
\section{Principe de routage par enveloppe inférieure}
% =================================================

\begin{definition}[Fonctionnelle de routage]
Pour une entrée \(w\in\Hh\), on introduit
\[
\mathcal E_{\mathrm{route}}(w;u,c):=\widetilde J_c(u;w),
\qquad (u,c)\in\Hh\times\Ccal.
\]
\end{definition}

\begin{theorem}[Principe de routage par enveloppe inférieure]
Pour tout \(w\in\Hh\), on a
\[
F(w):=\min_{c\in\Ccal}\widetilde V_c(w)
=
\inf_{(u,c)\in\Hh\times\Ccal}\mathcal E_{\mathrm{route}}(w;u,c).
\]
De plus, tout minimiseur global \((\hat u(w),\hat c(w))\) vérifie
\[
\hat u(w)=u_{\hat c(w)}^\star(w),
\qquad
\hat c(w)\in\argmin_{c\in\Ccal}\widetilde V_c(w).
\]
La fonction \(F\) est continue et localement lipschitzienne. Elle est directionnellement différentiable et, pour tout \(h\in\Hh\),
\[
F'(w;h)=\min_{c\in\Acal(w)}\langle \nabla \widetilde V_c(w),h\rangle.
\]
Enfin,
\[
\partial^C F(w)=\conv\{\nabla \widetilde V_c(w): c\in\Acal(w)\}.
\]
Sur toute cellule stricte \(\Omega_c^{\mathrm{str}}\), la fonction \(F\) est de classe \(C^1\) et
\[
\nabla F(w)=\nabla \widetilde V_c(w).
\]
\end{theorem}

\begin{proof}
Comme \(\Ccal\) est fini,
\[
\inf_{(u,c)\in\Hh\times\Ccal}\mathcal E_{\mathrm{route}}(w;u,c)
=
\min_{c\in\Ccal}\inf_{u\in\Hh}\widetilde J_c(u;w)
=
\min_{c\in\Ccal}\widetilde V_c(w).
\]
Le minimiseur en \(u\) est unique pour chaque contexte \(c\), d'où la première partie. Comme \(F\) est le minimum fini de fonctions \(C^1\), il est continu, localement lipschitzien et directionnellement différentiable. Les formules de dérivée directionnelle et de sous-différentiel de Clarke sont les règles standards du minimum fini de fonctions \(C^1\) localement lipschitziennes \cite{Clarke1983}. Sur une cellule stricte, une seule branche est active, donc \(F\) y coïncide avec une fonction \(C^1\) unique.
\end{proof}

\begin{remark}
Le résultat précédent identifie un routage exact à partir d'une lower envelope. La nouveauté ne réside pas dans l'égalité \(\inf_{u,c}=\min_c\inf_u\), mais dans son interprétation comme règle de décision et comme lecture variationnelle du noyau PCH.
\end{remark}

\begin{proposition}[Certificat de sélection par bornes]
Supposons que, pour chaque contexte \(c\), on dispose de bornes \(L_c(w)\le \widetilde V_c(w)\le U_c(w)\). S'il existe un indice \(c^\star\) tel que
\[
U_{c^\star}(w) < \min_{c\neq c^\star} L_c(w),
\]
alors \(c^\star\) est l'unique contexte routé en \(w\).
\end{proposition}

\begin{proof}
Si \(U_{c^\star}(w) < L_c(w)\) pour tout \(c\neq c^\star\), alors
\[
\widetilde V_{c^\star}(w)\le U_{c^\star}(w)<L_c(w)\le \widetilde V_c(w),
\]
ce qui implique que \(c^\star\) minimise strictement toutes les autres valeurs.
\end{proof}

% =================================================
\section{Géométrie des cellules, stabilité et robustesse}
% =================================================

\begin{proposition}[Correspondance active et version finie de Berge]
La correspondance active \(\Acal:\Hh\rightrightarrows\Ccal\) a des valeurs non vides et finies. Elle est semi-continue supérieurement. En particulier, sur toute cellule de dominance stricte, elle est localement constante. Ce point est la version finie du théorème du maximum de Berge \cite{Berge1963}.
\end{proposition}

\begin{proof}
Pour chaque \(c\),
\[
\{w\in\Hh:\ c\in\Acal(w)\}
=\bigcap_{j\in\Ccal}\{w\in\Hh:\widetilde V_c(w)\le \widetilde V_j(w)\},
\]
intersection finie d'ensembles fermés, donc ensemble fermé. Comme \(\Ccal\) est fini, la correspondance est à valeurs finies, non vides et fermées ; elle est donc semi-continue supérieurement. Sur une cellule stricte, les inégalités sont strictes et stables par continuité, d'où la constance locale.
\end{proof}

\begin{definition}[Cellules de contexte]
On définit
\[
\Omega_c^{\mathrm{str}}:=\{w\in\Hh:\ \widetilde V_c(w)<\widetilde V_j(w)\ \forall j\neq c\},
\]
\[
\Omega_c:=\{w\in\Hh:\ c\in\argmin_{j\in\Ccal}\widetilde V_j(w)\}
=\bigcap_{j\neq c}\{w\in\Hh:\widetilde V_c(w)\le \widetilde V_j(w)\},
\]
et
\[
\Bcal:=\{w\in\Hh:\ \#\Acal(w)\ge 2\}.
\]
\end{definition}

\begin{proposition}[Partition de l'espace latent et frontières]
Pour chaque \(c\in\Ccal\), \(\Omega_c^{\mathrm{str}}\) est ouvert et \(\Omega_c\) est fermé. Les cellules strictes sont deux à deux disjointes, et
\[
\Hh=\left(\bigsqcup_{c\in\Ccal}\Omega_c^{\mathrm{str}}\right)\sqcup \Bcal.
\]
Pour chaque paire \(c\neq j\), en posant
\[
\Gamma_{c,j}(w):=\widetilde V_c(w)-\widetilde V_j(w),
\qquad
\Sigma_{c,j}:=\{w\in\Hh:\Gamma_{c,j}(w)=0\},
\]
on a que \(\Sigma_{c,j}\) est fermé. Si \(\nabla\Gamma_{c,j}(w_0)\neq 0\), alors \(\Sigma_{c,j}\) est localement une hypersurface \(C^1\) au voisinage de \(w_0\).
\end{proposition}

\begin{proof}
Chaque \(\widetilde V_c-\widetilde V_j\) est continue. Donc \(\Omega_c^{\mathrm{str}}\), défini par des inégalités strictes de fonctions continues, est ouvert, et \(\Omega_c\), défini par des inégalités larges, est fermé. Les cellules strictes sont disjointes par contradiction immédiate. La fermeture de \(\Sigma_{c,j}\) découle de la continuité de \(\Gamma_{c,j}\). Si \(\nabla\Gamma_{c,j}(w_0)\neq 0\), le théorème des fonctions implicites s'applique à \(\Gamma_{c,j}(w)=0\).
\end{proof}

\begin{corollary}[Stabilité sous marge stricte]
Supposons qu'en un point \(w\in\Hh\), le minimiseur de \(c\mapsto\widetilde V_c(w)\) soit unique, noté \(c^\star\), et qu'il existe une marge
\[
\delta(w):=\min_{c\neq c^\star}\bigl(\widetilde V_c(w)-\widetilde V_{c^\star}(w)\bigr)>0.
\]
Alors il existe un voisinage \(U\) de \(w\) tel que tout \(z\in U\) admet le même contexte routé \(c^\star\).
\end{corollary}

\begin{proof}
Par continuité finie des fonctions \(\widetilde V_c\), on choisit \(U\) tel que \(|\widetilde V_c(z)-\widetilde V_c(w)|<\delta(w)/3\) pour tous \(c\) et tout \(z\in U\). L'argument standard de comparaison des branches donne alors la stabilité locale du contexte optimal.
\end{proof}

\begin{theorem}[Robustesse quantitative du routage]
Supposons qu'il existe, dans un voisinage de \(w\), une constante \(L<\infty\) telle que chaque fonction \(\widetilde V_c\) soit \(L\)-Lipschitz. Si le minimiseur de \(c\mapsto\widetilde V_c(w)\) est unique, noté \(c^\star\), et si sa marge vérifie \(\delta(w)>0\), alors tout point \(z\) tel que
\[
\|z-w\|<\frac{\delta(w)}{2L}
\]
admet le même contexte routé \(c^\star\).
\end{theorem}

\begin{proof}
Pour tout \(c\), on a \(|\widetilde V_c(z)-\widetilde V_c(w)|\le L\|z-w\|<\delta(w)/2\). Ainsi,
\[
\widetilde V_c(z)\ge \widetilde V_c(w)-\delta(w)/2\ge \widetilde V_{c^\star}(w)+\delta(w)-\delta(w)/2\ge \widetilde V_{c^\star}(z),
\]
pour tout \(c\neq c^\star\), ce qui montre la stabilité du contexte routé.
\end{proof}

\begin{theorem}[Optimalité décisionnelle du routage variationnel]
Supposons qu'il existe un couple aléatoire \((Z,W)\) et des pertes mesurables \(\ell_c\) telles que
\[
\widetilde V_c(w)=\mathbb E\bigl[\ell_c(Z,W)\mid W=w\bigr]
\]
pour tout contexte \(c\). Alors, pour tout sélecteur déterministe mesurable \(s:\Hh\to\Ccal\), on a
\[
\mathbb E\bigl[\widetilde V_{s(W)}(W)\bigr]\ge \mathbb E\bigl[F(W)\bigr].
\]
De plus, l'égalité a lieu si et seulement si \(s(W)\in\argmin_{c\in\Ccal}\widetilde V_c(W)\) presque sûrement. Ainsi, le routage dur par \(\argmin\) est Bayes-optimal parmi les sélecteurs déterministes au sens du risque conditionnel.
\end{theorem}

\begin{proof}
Pour presque tout \(w\), on a par définition
\[
\widetilde V_{s(w)}(w)\ge \min_{c\in\Ccal}\widetilde V_c(w)=F(w).
\]
En intégrant cette inégalité, on obtient le résultat. L'égalité pointwise exige que \(s(w)\) appartienne à l'ensemble des minimisateurs presque partout.
\end{proof}

\begin{corollary}[Robustesse sous perturbation des valeurs]
Supposons que des valeurs approchées \(\widehat V_c\) satisfassent, en un point \(w\),
\[
\sup_{c\in\Ccal}|\widehat V_c(w)-\widetilde V_c(w)|\le \varepsilon.
\]
Si le minimiseur \(c^\star\) est unique et si \(\delta(w)>2\varepsilon\), alors
\[
\argmin_{c\in\Ccal}\widehat V_c(w)=\{c^\star\}.
\]
\end{corollary}

\begin{proof}
Pour tout \(c\neq c^\star\),
\[
\widehat V_c(w)-\widehat V_{c^\star}(w)
\ge \widetilde V_c(w)-\widetilde V_{c^\star}(w)-2\varepsilon
\ge \delta(w)-2\varepsilon>0.
\]
Donc le minimiseur ne change pas.
\end{proof}

% =================================================
\section{Routage souple et complexité de calcul}
% =================================================

Le routage dur est discret. Pour obtenir une relaxation convexe, on passe au simplexe des contextes. Le routage souple est une approximation stable du routage exact, utile quand l'évaluation exhaustive de toutes les branches est coûteuse.

\begin{definition}[Simplexe des contextes]
On note
\[
\Delta_K:=\left\{p\in\R^K:\ p_c\ge 0,\ \sum_{c=1}^K p_c=1\right\},
\qquad
\Delta_K^\circ:=\{p\in\Delta_K:\ p_c>0\ \text{pour tout }c\}.
\]
\end{definition}

\begin{definition}[Fonctionnelle de routage souple]
Pour \(\beta>0\), on définit
\[
\Rcal_\beta(w,p)
=
\sum_{c=1}^K p_c\,\widetilde V_c(w)
+\frac{1}{\beta}\sum_{c=1}^K p_c\log p_c,
\qquad p\in\Delta_K,
\]
avec la convention \(0\log 0=0\).
\end{definition}

\begin{proposition}[Existence et unicité du routage souple]
Pour tout \(w\), la fonctionnelle \(p\mapsto \Rcal_\beta(w,p)\) admet un unique minimiseur \(p^\star_\beta(w)\) dans \(\Delta_K\).
\end{proposition}

\begin{proof}
La fonction \(p\mapsto \sum_c p_c\log p_c\) est strictement convexe sur \(\Delta_K\) et \(p\mapsto \sum_c p_c\,\widetilde V_c(w)\) est affine. La somme est donc strictement convexe sur le compact \(\Delta_K\), d'où l'existence et l'unicité du minimiseur.

Montrons maintenant que le minimiseur appartient à l'intérieur \(\Delta_K^\circ\). Supposons au contraire que \(p^\star_i=0\) pour un certain indice \(i\). Comme \(\sum_c p^\star_c=1\), il existe un indice \(j\) tel que \(p^\star_j>0\). Pour \(t\in(0,p^\star_j)\), posons \(p(t)=p^\star+t(e_i-e_j)\in\Delta_K\). On obtient une fonction unidimensionnelle dont la dérivée à droite en \(0\) est \(-\infty\), car \(\log t\to-\infty\) lorsque \(t\downarrow 0\). Cela contredit la minimalité de \(p^\star\). Ainsi, tout minimiseur appartient à \(\Delta_K^\circ\).
\end{proof}

\begin{theorem}[Forme explicite du routage souple]
Le minimiseur unique de \(\Rcal_\beta(w,\cdot)\) est donné par
\[
 p_c^\star(w)=\frac{\exp\!\bigl(-\beta\,\widetilde V_c(w)\bigr)}
 {\sum_{j=1}^K \exp\!\bigl(-\beta\,\widetilde V_j(w)\bigr)}.
\]
\end{theorem}

\begin{proof}
Comme le minimiseur appartient à \(\Delta_K^\circ\), on peut appliquer la méthode des multiplicateurs de Lagrange au problème
\[
\min_{p\in\Delta_K^\circ}\ \sum_{c=1}^K p_c\,\widetilde V_c(w)+\frac{1}{\beta}\sum_{c=1}^K p_c\log p_c.
\]
On obtient, pour chaque \(c\),
\[
\widetilde V_c(w)+\frac{1}{\beta}(1+\log p_c)+\lambda=0,
\]
donc
\[
 p_c=\exp(-1-\beta\lambda)\,\exp\!\bigl(-\beta\widetilde V_c(w)\bigr).
\]
La contrainte \(\sum_c p_c=1\) donne la formule annoncée. Le résultat est la forme usuelle du softmin entropique \cite{BauschkeCombettes2017,RockafellarWets1998}.
\end{proof}

\begin{corollary}[Valeur du problème de routage et limite dure]
On a l'identité
\[
\inf_{p\in\Delta_K}\Rcal_\beta(w,p)
=
-\frac{1}{\beta}\log\!\left(\sum_{c=1}^K \exp\!\bigl(-\beta\,\widetilde V_c(w)\bigr)\right).
\]
De plus,
\[
F(w)\le -\frac{1}{\beta}\log\!\left(\sum_{c=1}^K e^{-\beta \widetilde V_c(w)}\right)\le F(w)+\frac{\log K}{\beta}.
\]
Enfin, si le minimiseur de \(c\mapsto \widetilde V_c(w)\) est unique, noté \(c^\star(w)\), alors
\[
 p^\star_\beta(w)\xrightarrow[\beta\to+\infty]{}e_{c^\star(w)}.
\]
\end{corollary}

\begin{proof}
C'est la valeur classique de l'entropie log-sum-exp. La borne sandwich est la propriété standard du \(\log\)-sum-exp. Pour la limite dure, on compare chaque poids non optimal au poids optimal ; les différences strictement positives dans les valeurs réduites sont écrasées exponentiellement lorsque \(\beta\to+\infty\).
\end{proof}

\begin{proposition}[Profil de calcul et comparaison conceptuelle]
Si le coût d'une résolution locale du contexte \(c\) à précision \(\varepsilon\) est noté \(T_c(\varepsilon)\), alors un routage exact par comparaison exhaustive coûte
\[
C_{\mathrm{exact}}(\varepsilon)=\sum_{c\in\Ccal}T_c(\varepsilon).
\]
Un routage appris par gate a un coût d'inférence approximatif
\[
C_{\mathrm{gate}}\approx T_{\mathrm{gate}}+T_{c^\star},
\]
mais il exige une phase d'apprentissage séparée et n'offre pas, en lui-même, de certificat de marge sur le contexte choisi. Le routage variationnel exact est donc plus coûteux à l'inférence, mais il est interprétable, certifiable par marge et n'introduit pas de module externe de décision.
\end{proposition}

\begin{proof}
Le premier terme suit de la comparaison explicite de tous les contextes. Le second profil reflète le schéma classique gate-plus-expert : un gate est évalué puis un expert est exécuté. La différence conceptuelle est que le gate appris doit être entraîné et peut se tromper hors distribution, alors que le routage variationnel compare directement les valeurs locales et peut être validé par marge.
\end{proof}

\begin{corollary}[Routage approché par certificat court]
On peut définir un routage approché par
\[
\hat c_N(w)\in\argmin_{c\in\Ccal}\widetilde V_c^{(N)}(w).
\]
Lorsque \(N\) est suffisamment grand, \(\hat c_N(w)\) coïncide avec le routage exact sous l'hypothèse d'unicité du minimiseur.
\end{corollary}

\begin{proof}
C'est une conséquence directe de la convergence des valeurs approchées vers les valeurs exactes et de la marge stricte dès que le contexte optimal est séparé des autres.
\end{proof}

\begin{quote}
\small
\textbf{Algorithme de routage calculable.}

Entrée : \(w\), ensemble fini \(\Ccal\), tolérance \(\varepsilon\).

1. Pour chaque contexte \(c\in\Ccal\), résoudre le problème local et calculer une approximation \(\widetilde V_c^{(N)}(w)\).

2. Choisir \(c^\star\in\argmin_{c\in\Ccal}\widetilde V_c^{(N)}(w)\).

3. Si une marge \(\delta>0\) est certifiée, arrêter dès que
\[
\bigl|\widetilde V_c^{(N)}(w)-\widetilde V_{c^\star}^{(N)}(w)\bigr|<\delta/3
\qquad\text{pour tout }c\neq c^\star.
\]

Sortie : contexte routé \(c^\star\) et minimiseur local \(u_{c^\star}^\star(w)\).
\end{quote}

\begin{remark}
La procédure précédente n'est pas une heuristique cachée : elle n'utilise que des valeurs locales, une marge calculable et un critère d'arrêt explicite.
\end{remark}

% =================================================
\section{Pourquoi cette formulation n'est pas un simple \emph{Mixture of Experts}}
% =================================================

\begin{table}[ht]
\centering
\begin{tabular}{@{}lll@{}}
\toprule
Aspect & \emph{Mixture of Experts} & Routage PCH \\
\midrule
Sélection & gate appris & \(\argmin_c \widetilde V_c(w)\) \\
Objet de base & score séparé & fonction valeur locale \\
Décision & statistique / discriminative & variationnelle / énergétique \\
Interprétation & souvent opaque & géométrie de cellules \\
Certificat & rarement explicite & marge énergétique explicite \\
Apprentissage du routeur & requis & non requis pour la règle de base \\
\bottomrule
\end{tabular}
\caption{Différence conceptuelle entre routage appris et routage variationnel.}
\end{table}

Le message n'est pas qu'un gate appris serait inutile dans toute application, mais qu'il s'agit d'un objet différent : le routeur variationnel ne prédit pas le contexte, il le sélectionne en minimisant une énergie locale déjà calculée. La différence n'est pas seulement technique ; elle est paradigmatique. Ici, la décision est lue dans la structure variationnelle elle-même.

\begin{remark}
Le changement de paradigme se situe précisément à l'endroit suivant : dans un MoE, on apprend un routeur ; ici, on déduit une règle de sélection à partir d'une famille de fonctions valeur. Cette différence suffit à changer le statut mathématique de la décision.
\end{remark}

% =================================================
\section{PCH-ESM comme instanciation naturelle}
% =================================================

Dans PCH-ESM, une décomposition orthogonale et déterministe fournit une famille finie de sous-espaces spécialisés. Ce sont des contextes naturels pour le routage variationnel. Cette section doit être lue comme une instanciation de la théorie générale, non comme un second cœur du papier.

\begin{assumption}[Décomposition spectrale]
On suppose un espace latent \(\Hh=\Hh_0\oplus \Hh_1\oplus\cdots\oplus \Hh_n\), obtenu par projecteurs spectraux orthogonaux. On note \(\Pi_i:\Hh\to\Hh_i\) la projection orthogonale.
\end{assumption}

\begin{definition}[Bloc certifié]
On appelle bloc certifié un bloc \(i\) muni d'une valeur réduite \(\widetilde V_i^{\mathrm{ESM}}(w)\), d'une borne d'erreur explicite \(\epsilon_i(w)\), et d'une marge de séparation \(\delta_i(w)\) vérifiée sur le jeu de calibration.
\end{definition}

\begin{definition}[Contextes PCH-ESM]
On note \(\Ccal_{\mathrm{ESM}}=\{1,\dots,n\}\) l'ensemble des blocs spécialisés induits par la décomposition spectrale. À chaque bloc \(i\) on associe une valeur réduite \(\widetilde V_i^{\mathrm{ESM}}(w)\) calculée par le problème local du bloc correspondant.
\end{definition}

\begin{proposition}[Instantiation du routage sur PCH-ESM]
Si les contextes admissibles sont pris parmi les blocs spécialisés de PCH-ESM, alors le bloc routé est
\[
i^\star(w)\in\argmin_{i\in\Ccal_{\mathrm{ESM}}}\widetilde V_i^{\mathrm{ESM}}(w).
\]
Autrement dit, la décomposition spectrale fournit une famille de contextes admissibles, et le routage variationnel sélectionne le bloc de plus faible valeur réduite.
\end{proposition}

\begin{proof}
C'est l'application directe du principe de routage à la famille finie des blocs spécialisés.
\end{proof}

\begin{proposition}[Spécialisation orthogonale et gain de reconstruction]
Supposons que la cible vraie \(\theta^\star\) appartienne à un sous-espace spécialisé \(\Hh_i\), et qu'une estimation dense \(\hat\theta_{\mathrm{dense}}\in\Hh\) soit projetée par \(\hat\theta_{\mathrm{spec}}:=\Pi_i\hat\theta_{\mathrm{dense}}\). En notant \(\Pi_i^\perp:=\mathrm{Id}-\Pi_i\), on a alors
\[
\|\hat\theta_{\mathrm{dense}}-\theta^\star\|^2
=
\|\hat\theta_{\mathrm{spec}}-\theta^\star\|^2
+\|\Pi_i^\perp \hat\theta_{\mathrm{dense}}\|^2.
\]
En particulier, si \(\Pi_i^\perp \hat\theta_{\mathrm{dense}}\neq 0\), la spécialisation réduit strictement l'erreur quadratique.
\end{proposition}

\begin{proof}
Comme \(\theta^\star\in\Hh_i\), la différence \(\hat\theta_{\mathrm{dense}}-\theta^\star\) se décompose orthogonalement en une composante dans \(\Hh_i\) et une composante orthogonale. L'identité suit par Pythagore.
\end{proof}

\begin{proposition}[Collage pondéré inverse-variance]
Supposons que les certificats locaux de la composante partagée s'écrivent
\[
u_0^{(i)} = u_0^\star + \varepsilon_i,\qquad i=1,\dots,n,
\]
où \(u_0^\star\in\Hh_0\) est la vraie composante partagée, les erreurs \(\varepsilon_i\) sont indépendantes, centrées et isotropes, avec \(\mathbb E\|\varepsilon_i\|^2 = d\,\tau_i\) pour des \(\tau_i>0\). Alors, parmi tous les estimateurs affines non biaisés \(\tilde u_0=\sum_{i=1}^n \beta_i u_0^{(i)}\) avec \(\sum_i\beta_i=1\), le choix
\[
\beta_i^\star=\frac{\tau_i^{-1}}{\sum_{j=1}^n \tau_j^{-1}}
\]
minimise l'erreur quadratique moyenne. Le consensus pondéré est donc l'estimateur linéaire non biaisé optimal.
\end{proposition}

\begin{proof}
On a \(\mathbb E[\tilde u_0-u_0^\star]=0\) et, par indépendance,
\[
\mathbb E\|\tilde u_0-u_0^\star\|^2
=
 d\sum_{i=1}^n \beta_i^2 \tau_i.
\]
La minimisation sous la contrainte \(\sum_i\beta_i=1\) donne \(\beta_i\propto \tau_i^{-1}\), donc la formule annoncée.
\end{proof}

\begin{remark}
PCH-ESM est ici une instanciation structurée de la théorie générale : il ne remplace pas le cœur variationnel du manuscrit, il le spécialise sur des blocs spectraux certifiés.
\end{remark}

\begin{figure}[ht]
\centering
\begin{tikzpicture}[node distance=1.5cm, >=Latex, every node/.style={align=center}]
\node (pch) [draw, rounded corners, minimum width=3.1cm, minimum height=0.95cm] {PCH\\valeurs locales};
\node (esm) [draw, rounded corners, right=1.8cm of pch, minimum width=3.4cm, minimum height=0.95cm] {PCH-ESM\\blocs spécialisés};
\node (vals) [draw, rounded corners, right=1.8cm of esm, minimum width=3.4cm, minimum height=0.95cm] {Valeurs réduites\\par bloc};
\node (sel) [draw, rounded corners, below=1.2cm of vals, minimum width=3.4cm, minimum height=0.95cm] {Sélection\\\(\argmin\)};
\node (out) [draw, rounded corners, left=1.8cm of sel, minimum width=3.4cm, minimum height=0.95cm] {Contexte choisi\\ou bloc expert};

\draw[->] (pch) -- (esm);
\draw[->] (esm) -- (vals);
\draw[->] (vals) -- (sel);
\draw[->] (sel) -- (out);
\end{tikzpicture}
\caption{Chaîne conceptuelle PCH \(\rightarrow\) PCH-ESM \(\rightarrow\) routage.}
\end{figure}

% =================================================
\section{Exemples}
% =================================================

\begin{example}[Cas unidimensionnel symétrique]
On prend \(d=1\), \(\Mcal_1=\{-a\}\), \(\Mcal_2=\{a\}\) avec \(a>0\), et des paramètres identiques
\[
\lambda_{\mathrm{ctx},1}=\lambda_{\mathrm{ctx},2}=\lambda>0,
\qquad
\lambda_{\mathrm{ridge},1}=\lambda_{\mathrm{ridge},2}=\rho\ge 0.
\]
Alors \(\mu_1=\mu_2=1+2\rho\), et une complétion du carré donne
\[
u_1^\star(w)=\frac{w-2\lambda a}{1+2\lambda+2\rho},
\qquad
u_2^\star(w)=\frac{w+2\lambda a}{1+2\lambda+2\rho}.
\]
De plus,
\[
\widetilde V_1(w)-\widetilde V_2(w)
=
\frac{4a\lambda}{(1+2\rho)(1+2\lambda+2\rho)}\,w.
\]
Par conséquent,
\[
\Omega_1^{\mathrm{str}}=(-\infty,0),
\qquad
\Omega_2^{\mathrm{str}}=(0,\infty),
\qquad
\Bcal=\{0\}.
\]
Le routage exact se lit donc immédiatement : le contexte \(1\) est choisi pour \(w<0\), le contexte \(2\) pour \(w>0\), et la frontière de décision est réduite au point \(0\).
\end{example}

\begin{example}[Cas bidimensionnel à trois contextes]
On prend \(\Hh=\R^2\) et trois fonctions valeur convexes
\[
\widetilde V_1(x_1,x_2)=\|(x_1,x_2)-(1,0)\|^2,
\qquad
\widetilde V_2(x_1,x_2)=\|(x_1,x_2)-(-1,0)\|^2,
\]
\[
\widetilde V_3(x_1,x_2)=\|(x_1,x_2)-(0,1)\|^2.
\]
Alors les frontières sont données par des égalités affines entre ces trois branches ; la partition comprend des cellules polygonales dont la forme n'est plus simplement symétrique. Ce toy model montre que la géométrie de routage ne se limite pas à un seul hyperplan de séparation.
\end{example}

\begin{figure}[ht]
\centering
\begin{tikzpicture}[scale=1]
\draw[->] (-3.2,0) -- (3.2,0) node[right] {$x_1$};
\draw[->] (0,-2.0) -- (0,2.2) node[above] {$x_2$};
\fill[blue!10] (-3,-2) rectangle (0,2);
\fill[orange!12] (0,-2) rectangle (3,2);
\fill[green!10] (-3,-2) -- (0,2) -- (3,-2) -- cycle;
\draw[dashed, very thick] (0,-2) -- (0,2);
\draw[thick] (-1,0.5) -- (0.5,1.5) -- (1.5,0.2);
\node at (-1.4,1.6) {$\Omega_1^{\mathrm{str}}$};
\node at (1.35,1.6) {$\Omega_2^{\mathrm{str}}$};
\node at (-0.6,-0.2) {$\Omega_3^{\mathrm{str}}$};
\node at (0.35,0.2) {$\Bcal$};
\end{tikzpicture}
\caption{Exemple bidimensionnel : la géométrie de routage n'est pas limitée aux cas symétriques.}
\end{figure}

% =================================================
\section{Relation à la littérature}
% =================================================

Le cadre ici présenté se place à l'interface de l'analyse convexe, de la théorie des fonctions valeur, de la programmation disjonctive et des modèles de type \emph{Mixture of Experts}. Les références classiques pertinentes sont celles de Moreau, Rockafellar, Berge, Clarke, ainsi que Bauschke--Combettes et Rockafellar--Wets \cite{Moreau1965,Rockafellar1970,Berge1963,Clarke1983,BauschkeCombettes2017,RockafellarWets1998}.

La nouveauté du manuscrit n'est pas d'inventer de nouveaux outils de convexité ; elle est de faire du routage un problème variationnel à part entière. Les méthodes classiques de routage utilisent un score appris, un classifieur ou un \emph{gate} séparé. Ici, le contexte est sélectionné comme minimiseur d'une famille de valeurs réduites issues directement de problèmes locaux. Le routage n'est donc pas un module additionnel de décision ; il est la lecture variationnelle de cette famille de valeurs.

En termes pratiques, le routage appris est souvent moins coûteux à l'inférence une fois le gate entraîné, mais il introduit une phase d'apprentissage séparée, une dépendance aux données d'entraînement du gate et une opacité de décision relative. Le routage variationnel est plus coûteux si l'on résout exhaustivement toutes les branches, mais il est interprétable, il s'inscrit directement dans PCH et il peut être accéléré par un pré-roulage court avec marge. Ce compromis est la formulation scientifique du problème.

\begin{remark}
Cette lecture permet de séparer proprement ce qui relève d'un rappel classique et ce qui relève de la contribution propre du manuscrit : la sélection de contexte par minimisation des valeurs réduites et l'interprétation géométrique qui en découle.
\end{remark}

% =================================================
\section{Conclusion}
% =================================================

Le noyau PCH n'est pas le problème à résoudre ; il est la brique locale qui fournit des fonctions valeur comparables. Le routage est la continuité naturelle de cette structure : le contexte sélectionné est celui qui minimise la valeur locale réduite.

La formulation canonique est
\[
(\hat u(w),\hat c(w))\in\argmin_{(u,c)\in\Hh\times\Ccal}\mathcal E_{\mathrm{route}}(w;u,c),
\]
équivalente à
\[
\hat c(w)\in\argmin_{c\in\Ccal}\widetilde V_c(w),
\qquad
\widetilde V_c(w)=e_{\mu_c g_c}(w).
\]
La version souple est donnée par
\[
p^\star_\beta(w)=\argmin_{p\in\Delta_K}\Rcal_\beta(w,p).
\]

Cette formulation conserve strictement la normalisation, la lower envelope, la géométrie de cellules, la relaxation convexe sur le simplexe et l'instanciation PCH-ESM. Le message final est mesuré : la nouveauté ne vient pas des théorèmes classiques eux-mêmes, mais de leur assemblage pour faire émerger un routage variationnel à partir des fonctions valeur locales.

\begin{thebibliography}{99}

\bibitem{Moreau1965}
J.-J. Moreau,
\newblock Proximité et dualité dans un espace hilbertien.
\newblock \emph{Bull. Soc. Math. France}, 93:273--299, 1965.

\bibitem{Rockafellar1970}
R. T. Rockafellar,
\newblock \emph{Convex Analysis}.
\newblock Princeton University Press, 1970.

\bibitem{Berge1963}
C. Berge,
\newblock \emph{Topological Spaces: Including a Treatment of Multi-Valued Functions, Vector Spaces and Convexity}.
\newblock Oliver \& Boyd, 1963.

\bibitem{Clarke1983}
F. H. Clarke,
\newblock \emph{Optimization and Nonsmooth Analysis}.
\newblock Wiley, 1983.

\bibitem{BauschkeCombettes2017}
H. H. Bauschke and P. L. Combettes,
\newblock \emph{Convex Analysis and Monotone Operator Theory in Hilbert Spaces}.
\newblock 2nd ed., Springer, 2017.

\bibitem{RockafellarWets1998}
R. T. Rockafellar and R. J.-B. Wets,
\newblock \emph{Variational Analysis}.
\newblock Springer, 1998.

\end{thebibliography}

\end{document}
